
\chapter*{EXPOSÉ XIX\\[0.4em]
REDUCTIVE GROUPS---GENERALITIES}
\addcontentsline{toc}{chapter}{Exposé XIX. Reductive Groups---Generalities}

\begin{center}
\textit{by M. Demazure}
\end{center}

\begingroup
\renewcommand{\thefootnote}{0}%
\footnotetext{Editor's note: version of 13 October 2024.}%
\endgroup

The remainder of this Seminar is devoted to the study of reductive groups.  Its
principal aim is to generalize Chevalley's classical results (the
\emph{Bible} and \emph{Tôhoku})%
\footnote{Editor's note: see the bibliographical references at the end of this
Exposé.  In particular, the 2005 re-edition of the Chevalley Seminar
1956--1958, cited as \emph{Bible}, was revised by P.\ Cartier; see [Ch05].}
to arbitrary base schemes.  The two central results are the uniqueness theorem
(Exposé XXIII, Theorem 4.1 and \SGARefLink{sga3:xix:paragraph:5-1}{Corollaries 5.1}--5.10) and the existence theorem
(\SGARefLink{sga3:xxv:theorem:1:1}{Exposé XXV, Theorem 1.1}) for pinned reductive group schemes corresponding to
prescribed root data.  The proofs used here are inspired by Chevalley's; the
language of schemes makes it possible to give them greater effectiveness.

The results of the first volume of the \emph{Bible} (Exposés 1--13) will be
used systematically.  By contrast, we shall prove directly over an arbitrary
scheme the results of the second volume (in particular the isomorphism
theorem); knowledge of the proofs over an algebraically closed field is
therefore not absolutely indispensable.

In proving these two fundamental results, we shall use only the most elementary
results from the theory of groups of multiplicative type, found essentially in
Exposés VIII and IX; on the other hand, we shall make essential use of the
results of Exposé XVIII.%
\footnote{Editor's note: more precisely, Proposition \SGARefLink{sga3:xviii:prop:2.3}{XVIII 2.3} (extension of a
``generic homomorphism'' between groups) is used in \SGARefLink{sga3:xxii:prop:4.1.11}{XXII 4.1.11} and then in the
proof of uniqueness in XXIII, as well as in XXIV.  Theorem \SGARefLink{sga3:xviii:thm:3.7}{XVIII 3.7}
(construction of a group from a group germ) is used only in XXV.}
The reader interested specifically in the existence and uniqueness theorems
may, on a first reading, skip Exposés X--XVII.

\SGASection{1}{Review of groups over an algebraically closed field}
\SGARefTarget{sga3:xix:section:1}

\noindent\SGAInlineRefTarget{sga3:xix:paragraph:1-1}\textbf{1.1.}\quad
Throughout this subsection, \(k\) will denote an algebraically closed field.
As announced above, the only results from the \emph{Bible} used below are
contained in its first volume (Exposés 1--13).
